\section*{Capítulo \ref{ch:g3:bones-and-muscles} --- Ossos e músculos}

\begin{solution}{exo:g3:bones-and-muscles:1}
O \omterm{def:g3:bones-and-muscles:skeleton}{crânio} protege o cérebro; a \omterm{def:g3:bones-and-muscles:skeleton}{coluna} segura o tronco em pé (uma
torre flexível); as \omterm{def:g3:bones-and-muscles:skeleton}{costelas} formam uma caixa de proteção em volta
do coração e dos pulmões; os \omterm{def:g3:bones-and-muscles:skeleton}{ossos} dos \omterm{def:g1:my-body:parts}{membros} sustentam os \omterm{def:g1:my-body:parts}{braços} e
as \omterm{def:g1:my-body:parts}{pernas}.
\end{solution}

\begin{solution}{exo:g3:bones-and-muscles:2}
Cerca de $200$. O mais comprido é o \omterm{def:g3:bones-and-muscles:skeleton}{osso} da coxa.
\end{solution}

\begin{solution}{exo:g3:bones-and-muscles:3}
Não --- um \omterm{def:g3:bones-and-muscles:skeleton}{osso} dobrado é um \omterm{def:g3:bones-and-muscles:skeleton}{osso} quebrado. Todo o dobrar acontece
nas \omterm{prop:g3:bones-and-muscles:joints}{articulações}, onde dois \omterm{def:g3:bones-and-muscles:skeleton}{ossos} se encontram.
\end{solution}

\begin{solution}{exo:g3:bones-and-muscles:4}
Ela é uma torre de mais de $20$ ossinhos, cada um capaz de girar
alguns graus sobre os vizinhos; as muitas dobras pequenas somam uma
grande curva flexível.
\end{solution}

\begin{solution}{exo:g3:bones-and-muscles:5}
Ele encurta e engrossa, e assim puxa o \omterm{def:g3:bones-and-muscles:skeleton}{osso} em que está preso.
Nenhum \omterm{def:g3:bones-and-muscles:muscle}{músculo} consegue empurrar.
\end{solution}

\begin{solution}{exo:g3:bones-and-muscles:6}
Porque um \omterm{def:g3:bones-and-muscles:muscle}{músculo} só puxa, e portanto um \omterm{def:g3:bones-and-muscles:muscle}{músculo} sozinho só
consegue mover a \omterm{def:g1:my-body:joint}{articulação} para um lado. Ao dobrar: o \omterm{def:g3:bones-and-muscles:muscle}{músculo} da
frente se contrai e puxa o antebraço para cima, enquanto o de trás
relaxa. Ao esticar: o \omterm{def:g3:bones-and-muscles:muscle}{músculo} de trás se contrai enquanto o da
frente relaxa.
\end{solution}

\begin{solution}{exo:g3:bones-and-muscles:7}
O próprio \omterm{def:g3:bones-and-muscles:skeleton}{osso} vivo reconstrói a fratura devagar; o gesso só segura
os pedaços parados enquanto o \omterm{def:g3:bones-and-muscles:skeleton}{osso} trabalha. Isso mostra que o \omterm{def:g3:bones-and-muscles:skeleton}{osso}
é material vivo, alimentado pelo sangue --- e não morto como o giz.
\end{solution}

\begin{solution}{exo:g3:bones-and-muscles:8}
Ao dobrar: a frente do \omterm{def:g1:my-body:parts}{braço} incha e endurece (esse \omterm{def:g3:bones-and-muscles:muscle}{músculo} está se
contraindo). Ao esticar: a frente amolece enquanto a parte de trás
do \omterm{def:g1:my-body:parts}{braço} se aperta --- o par trocando de papel.
\end{solution}

\begin{solution}{exo:g3:bones-and-muscles:9}
Os \omterm{def:g1:my-body:joint}{joelhos} são \omterm{prop:g3:bones-and-muscles:joints}{articulações} feitas para dobras grandes, movidas
pelos \omterm{def:g3:bones-and-muscles:muscle}{músculos} mais fortes do corpo; a \omterm{def:g3:bones-and-muscles:skeleton}{coluna} é feita para dobras
suaves espalhadas por muitos ossinhos. Levantar com os \omterm{def:g1:my-body:joint}{joelhos}
dobrados e as costas retas dá o trabalho pesado às \omterm{def:g1:my-body:parts}{pernas} e poupa a
\omterm{def:g3:bones-and-muscles:skeleton}{coluna}.
\end{solution}

\begin{solution}{exo:g3:bones-and-muscles:10}
Cada fio, como cada \omterm{def:g3:bones-and-muscles:muscle}{músculo}, só consegue puxar. A resposta de quem
fabrica marionetes é a mesma do biólogo: dois fios por movimento ---
um para puxar o \omterm{def:g1:my-body:parts}{braço} para cima, outro para puxar de volta para
baixo --- trabalhando por turnos, exatamente como os pares de
\omterm{def:g3:bones-and-muscles:muscle}{músculos} do corpo.
\end{solution}
